% !PRINTED 197
\weberpage{197}

The function (1) vanishes for all values $u^0$ of $u$ different from zero
which satisfy the condition
\[
\pm n\,u^0 \equiv u^0 \pmod{\omega_1,\omega_2}
\]
hence for
\weq{4}{u^0 = \frac{\nu_1\omega_1 + \nu_2\omega_2}{n\pm 1},}
where again $\nu_1, \nu_2$ are arbitrary integers. The zero points are of
the first order, the points of infinity of the second.

We therefore now introduce a function $\psi_n(u)$, which we define in the
following way:
\weq{5}{\psi_n(u)^2 = n^2 \prod^{\nu_1,\nu_2}\left[\wp(u) - \wp\left(\frac{\nu_1\omega_1+\nu_2\omega_2}{n}\right)\right],}
where $\nu_1, \nu_2$ each run through a complete system of residues with
respect to the modulus $n$, with the sole exclusion of the pair of values
$0,0$; and we add the further stipulation that $\psi_1 = 1$ shall hold.

If $n$ is odd, then in the product (5) every linear factor occurs twice,
since the two incongruent values
\[
u = \pm\frac{\nu_1\omega_1 + \nu_2\omega_2}{n}
\]
give the same value of $\wp(u)$. But if $n$ is even, then again all linear
factors occur twice in (5), with the exception of the three
\[
\wp(u) - \wp\left(\frac{\omega_1}{2}\right),\quad
\wp(u) - \wp\left(\frac{\omega_2}{2}\right),\quad
\wp(u) - \wp\left(\frac{\omega_1+\omega_2}{2}\right),
\]
which occur only once. If one now observes that, according to § 41,
\[
\wp'(u)^2 = 4\left[\wp(u)-\wp\left(\frac{\omega_1}{2}\right)\right]\left[\wp(u)-\wp\left(\frac{\omega_2}{2}\right)\right]\left[\wp(u)-\wp\left(\frac{\omega_1+\omega_2}{2}\right)\right]
\]
holds, then the result follows:
\weq{6}{
\begin{array}{ll}
n\text{ odd:} & \psi_n = P_n,\\
n\text{ even:} & \psi_n = \wp'(u)\,P_n,
\end{array}}
where $P_n$ denotes a polynomial in $\wp(u)$ which for
odd $n$ is of degree $\frac12(n^2-1)$ and for even $n$ of degree
$\frac12(n^2-4)$, and in which the term of highest order is respectively
equal to
\[
n\,\wp(u)^{\frac{n^2-1}{2}},\qquad -\frac{n}{2}\,\wp(u)^{\frac{n^2-4}{2}}
\]
%CONT>>
